\section{The Treatise on Law, II: Eternal and Natural Law (ST I-II qq.~93--94)}

\subsection{Question 93: the eternal law as measure of all law}

Question 93 gives the eternal law its formal definition through the analogy of the craftsman. In every artificer there pre-exists the plan (\latin{ratio}) of what his art will make; in every governor, the plan of the order to be carried out by the governed. God is both: creator, whose wisdom's plan has the character of art or exemplar; and governor of every act and motion, whose wisdom's plan, \latin{moventis omnia ad debitum finem}, has the character of law. Hence: \latin{lex aeterna nihil aliud est quam ratio divinae sapientiae, secundum quod est directiva omnium actuum et motionum}---the eternal law is nothing else than the plan of divine wisdom as directive of all acts and movements.\endnote{I-II q.~93 a.~1 co. ``Type'' in the Dominican translation renders \latin{ratio}---plan, idea, intelligible pattern. Note the precise difference from the exemplar: the same divine wisdom is \emph{art} as creative of things, \emph{law} as directive of acts toward ends.}

Article 3 draws the consequence that turns a theological thesis into a jurisprudence: whether every law is derived from the eternal law. In ordered movers, the power of the second mover derives from the first; in ordered governors, the plan of government descends from the chief---as the plan of what is to be done in a city flows from king to administrators, or from architect to craftsmen. Since the eternal law is the plan of government in the supreme governor, \latin{necesse est quod omnes rationes gubernationis quae sunt in inferioribus gubernantibus, a lege aeterna deriventur}. And so, in words that decide in advance everything the treatise will say about unjust law: \latin{omnes leges, inquantum participant de ratione recta, intantum derivantur a lege aeterna}---all laws, \emph{insofar as they partake of right reason}, are to that extent derived from the eternal law; with Augustine, in temporal law nothing is just and lawful which men have not derived to themselves from the eternal law.\endnote{I-II q.~93 a.~3 co., quoting \emph{De libero arbitrio} I. The restriction \latin{inquantum participant de ratione recta} is load-bearing: derivation is not a historical pedigree but a present participation, and a positive enactment participates exactly as far as it is right. The same article's \latin{ad}~1 handles even the \latin{fomes}---the inclination to sin---as ``law'' only in the sense of a penalty deriving from divine justice, not a norm.}

Two boundary-markers should be posted here, because popular usage blurs them. First, the eternal law is not ``divine positive law writ large''; it is prior to the distinction between natural and revealed, and both are its participations---one through creation, one through history. Second, the eternal law is not available to us as a code we could consult; we reach it only in its participations. That is why the canonical tradition never cites ``eternal law'' as a working source of obligation the way it cites natural or divine positive law: the eternal law functions as the account of \emph{why} the others oblige, not as a fifth text.

\subsection{Question 94: the natural law in detail}

Question 94 is the treatise's most quoted question, and the article at its center (a.~2) repays exact reading, because what it actually says is more disciplined than what it is often made to say.

\textbf{The first precept and the order of precepts} (a.~2). The precepts of natural law stand to practical reason as the first principles of demonstration stand to speculative reason: both are \latin{per se nota}, self-evident. As ``being'' is the first thing apprehended absolutely, so ``good'' is the first thing apprehended by practical reason, which is ordered to action: every agent acts for an end under the aspect of good. Therefore the first principle of practical reason is founded on the notion of good---\latin{bonum est quod omnia appetunt}---and the first precept of law is: \latin{bonum est faciendum et prosequendum, et malum vitandum}---good is to be done and pursued, and evil avoided. All other precepts of the natural law are founded on this, so that whatever practical reason naturally apprehends as human good belongs to its precepts. And because good has the nature of an end, \latin{secundum ordinem inclinationum naturalium, est ordo praeceptorum legis naturae}: the order of the precepts follows the order of the natural inclinations---first what man shares with all substances (the conservation of his being), then what he shares with the animals (the union of male and female, the rearing of offspring), then what is proper to reason (to know the truth about God, to live in society).\endnote{I-II q.~94 a.~2 co. Two exact points. First, the first precept is not a platitude but a \emph{principle}: it does not by itself tell anyone what is good; it is the form of every practical judgment, as non-contradiction is the form of every theoretical one. Second, the ``inclinations'' are not desires certifying themselves: they enter the law only as \emph{naturally apprehended by reason} as goods---the inclination ordered by reason, not the inclination as such, is the norm. Both points matter to the twentieth-century disputes reported in section 12.}

\textbf{Universality and its limits} (a.~4). Is the natural law the same in all? As to the common first principles, \latin{est eadem apud omnes et secundum rectitudinem, et secundum notitiam}---the same for all both in rectitude and in knowledge. As to the proper conclusions drawn from them, it is the same for all \latin{ut in pluribus}---in the majority of cases---but can fail in particulars, both in rectitude (returning a deposit is right, but not to a man demanding weapons to attack the fatherland) and in knowledge, where passion, custom, or corrupt habit obscures the conclusion.\endnote{I-II q.~94 a.~4 co., with the deposit example. The article's realism about moral knowledge---citing, in its final lines, peoples among whom even robbery was not reputed wrong---is the treatise's own answer to the objection ``if natural law were real, everyone would agree about it.'' Disagreement about conclusions is predicted by the theory, not fatal to it.}

\textbf{Immutability} (a.~5). Can the natural law be changed? By way of \emph{addition}, freely: many things useful to human life have been added over and above it, \latin{tam per legem divinam, quam etiam per leges humanas}. By way of \emph{subtraction}: as to first principles, \latin{lex naturae est omnino immutabilis}---altogether unchangeable; as to secondary precepts, unchangeable \latin{ut in pluribus}, though in some particular and rarer cases observance can be legitimately blocked.\endnote{I-II q.~94 a.~5 co. The famous hard cases in the \latin{ad}~2 (the divine commands concerning Isaac, the Egyptians' goods, Hosea's marriage) are resolved not by suspending the natural law but by locating God's dominion inside its terms: death, property, and marriage are what they are \latin{secundum legem divinitus traditam}, so the Lord of life, goods, and covenants does not steal or commit adultery in disposing of them. One may accept or contest the exegesis; the structure of the claim is what matters for the taxonomy---immutability of principle, divine administration of the matter.}

\textbf{Indelibility} (a.~6). Can it be blotted out of the human heart? As to common principles, \latin{nullo modo potest a cordibus hominum deleri in universali}---in no way, in the abstract; though it is blotted out in a particular action when passion hinders the application of principle to case. As to secondary precepts, it \emph{can} be deleted from hearts---by evil persuasions, as errors occur even about necessary conclusions in speculative matters, or \latin{propter pravas consuetudines et habitus corruptos}, by depraved customs and corrupt habits, as among those who did not reckon robbery, or vices against nature, to be sins.\endnote{I-II q.~94 a.~6 co., citing Romans 1. Read together, aa.~4--6 constitute a doctrine with juridical teeth: the natural law's universality is a universality of \emph{principle and ordinary case}, compatible with real cultures really losing real conclusions---which is exactly why the tradition holds that promulgated law, divine and human, is \emph{needed} to guard what is in principle knowable (compare CCC 1960 and q.~94 a.~5 ad 1's ``the written law is said to be given for the correction of the natural law''; section 11).}

What emerges from question 94 is a natural law with a definite architecture: a single self-evident first precept; primary precepts tracking the ordered natural inclinations; secondary precepts related to the primary as proximate conclusions; and, descending further, determinations that are no longer natural law at all but positive law built upon it. Everything in the canonical sections of this article---why divine law cannot be dispensed, why custom cannot prevail against it, why the Church claims to \emph{declare} rather than \emph{make} the impediments of natural law---presupposes precisely this architecture, and not a vaguer notion of ``natural law'' as whatever morality one takes to be obvious.
